\begin{frame}{Struktura relacyjna}

\begin{columns}
\column{0.65\textwidth}

    Dla grafu \(H=(V,E)\) definiujemy problem CSP\((H)\):
    \begin{align*}
        & \text{WEJŚCIE: Graf } G \\
        & \text{PYTANIE: Czy istnieje homomorfizm z } G \text{ do } H?
    \end{align*}

    \vspace{1em}
    Równoważnie można zapytać, czy graf \(G\) jest \(H\)-kolorowalny.

\column{0.35\textwidth}
    \begin{figure}
        \centering
        \begin{tikzpicture}[
    scale=0.65,
    line width=1.2pt,
    rotate=-90
]

    % Graph
    \begin{scope}[minimum size=0.25cm]
        \node[circle, draw=black, fill=cfblue!80] (v1) at (-1,1) {};
        \node[circle, draw=black, fill=cfteagreen!80] (v2) at (1,1) {};
        \node[circle, draw=black, fill=cfred!80] (v3) at (-2.4,0) {};
        \node[circle, draw=black, fill=cfblue!80] (v4) at (1,-1) {};
        \node[circle, draw=black, fill=cfred!80] (v5) at (2.5,0) {};
        \node[circle, draw=black, fill=cfteagreen!80] (v6) at (-1,-1) {};
                
        \draw (v1) -- (v2);
        \draw (v1) -- (v3);
        \draw (v3) -- (v6);
        \draw (v1) -- (v6);
        \draw (v2) -- (v5);
        \draw (v2) -- (v4);
        \draw (v6) -- (v4);
    \end{scope}


    % Mapping
    \begin{scope}[minimum size=1.2cm]
        \node[circle, draw=cfgrey!80, fill=cfgrey!80] (t1) at (5,4) {};
        \node[circle, draw=cfgrey!80, fill=cfgrey!80] (t3) at (3,2.5) {};
        \node[circle, draw=cfgrey!80, fill=cfgrey!80] (t2) at (5,1) {};
    \end{scope}

    \begin{scope}[minimum size=0.05cm]
        \node[circle, draw=black, fill=cfblue!80] (v11) at (4.75,3.75) {};
        \node[circle, draw=black, fill=cfblue!80] (v44) at (5.25,4.25) {};
        
        \node[circle, draw=black, fill=cfred!80] (v55) at (3.1,2.9) {};
        \node[circle, draw=black, fill=cfred!80] (v33) at (2.85,2.1) {};
    
        \node[circle, draw=black, fill=cfteagreen!80] (v22) at (4.75,0.65) {};
        \node[circle, draw=black, fill=cfteagreen!80] (v66) at (5.35,1.25) {};
    \end{scope}

    \draw[draw=cfgrey, line width=6] (t1) -- (t2);
    \draw[draw=cfgrey, line width=6] (t1) -- (t3);
    \draw[draw=cfgrey, line width=6] (t2) -- (t3);

    \draw (v11) -- (v22);
    \draw (v11) -- (v33);
    \draw (v33) -- (v66);
    \draw (v11) -- (v66);
    \draw (v22) -- (v55);
    \draw (v22) -- (v44);
    \draw (v66) -- (v44);

\end{tikzpicture}

    \end{figure}
\end{columns}
\vspace{3em}

\end{frame}